\subsubsection{碰撞核谱系：从递推到非负整数 realization（可审计协议）}\label{subsubsec:pom-collision-kernel-family}
上文已经给出一条“从折叠纤维到碰撞谱”的统一入口：
对任意整数 \(q\ge 2\)，
\[
S_q(m):=\sum_{x\in X_m}d_m(x)^q
=\#\{(\omega^{(1)},\dots,\omega^{(q)})\in\Omega_m^q:\ \Fold_m(\omega^{(1)})=\cdots=\Fold_m(\omega^{(q)})\}.
\]
其中 \(q=2\) 是成对碰撞，\(q=3\) 是三元碰撞，依此类推；\(S_q(m)\) 因而描述的是一个“\(q\)-重纤维积”上的计数对象，而非对某个固定核的简单幂函数。
\par
在 \(q=2,3,4\) 的情形，本文已经构造出低维的非负整数碰撞核 \(A_q\)，并把其 Perron 根 \(r_q=\rho(A_q)\) 与投影熵率 \(h_q=\log(2^q/r_q)\) 纳入同一 \(\zeta\)/primitive/有限部接口（见 \S\ref{subsec:pom-s2}--\S\ref{subsec:pom-s4} 及表 \ref{tab:fold_collision_renyi_spectrum}）。
另一方面，对 \(q\ge 5\) 已经核验存在小阶整系数递推（表 \ref{tab:fold_collision_moment_recursions}），从而 \(r_q\) 与 \(h_q\) 可精确抽取，但递推本身并不自动给出“非负整数图核”的 realization。
这引出一个自然的升级研究线：把“已知递推”进一步提升为“可解释的有限状态碰撞核谱系”。

\begin{proposition}[碰撞矩谱的递推—谱半径等价：Rényi 指纹率由 Perron 根族统一封闭]\label{prop:pom-collision-renyi-perron-closure}
固定整数 \(q\ge 2\)。
设碰撞矩 \(S_q(m)=\sum_x d_m(x)^q\) 自某个 \(m\) 起满足一条最小阶常系数线性递推（等价于 Hankel 秩有限），并记其递推特征多项式的最大模根为 \(r_q\)（Perron 根）。
则：
\begin{enumerate}
  \item 存在某个最小维线性实现（例如命题 \ref{prop:pom-hankel-realization-protocol} 的 Hankel realization）使得
  \(
  S_q(m)=\Theta(r_q^m)
  \)，并且该 \(r_q\) 等于任意最小实现矩阵的谱半径（推论 \ref{cor:pom-hankel-free-energy-evaluator}）。
  \item 若进一步存在自然的拼接次乘性（例如 \(S_q(m+n)\le S_q(m)S_q(n)\)），则极限
  \[
  \boxed{\ r_q=\lim_{m\to\infty}S_q(m)^{1/m}\ }
  \]
  存在且与上述 Perron 根一致。
  \item 定义 Rényi 指纹率
  \[
  h_q:=\log\!\left(\frac{2^q}{r_q}\right),
  \qquad
  h_q^{\mathrm R}:=\frac{h_q}{q-1}.
  \]
  则一旦 \(h_q^{\mathrm R}\) 的极限口径成立（例如由上式直接定义于极限 \(r_q\)），
  由有限支持分布的 Rényi 单调性可知 \(q\mapsto h_q^{\mathrm R}\) 单调递减，并且存在端点极限
  \[
  h_\infty^{\mathrm R}:=\lim_{q\to\infty}h_q^{\mathrm R}
  =\log 2-\lim_{q\to\infty}\frac{1}{q}\log r_q.
  \]
  在 Fold 门上，该端点由命题 \ref{prop:pom-rq-qinfty-endpoint} 锁定为 \(h_\infty^{\mathrm R}=\log(2/\sqrt{\varphi})\)（亦见注 \ref{rem:pom-renyi-gap-to-hinfty}）。
\end{enumerate}
\end{proposition}

\begin{proposition}[从递推到最小线性实现：Hankel realization]\label{prop:pom-hankel-realization-protocol}
固定 \(q\ge 2\)，设序列 \(S_q(m)\) 满足某个最小阶的常系数线性递推（等价于其 Hankel 秩有限）。
令 \(a_n:=S_q(n+2)\)（\(n\ge 0\)），并令 \(d:=\mathrm{rank}(H)\) 为其 Hankel 秩。
则存在一个维数等于 \(d\) 的最小线性实现
\[
a_n=\alpha^{\mathsf T}A^n\beta\qquad(n\ge 0),
\]
其中 \(A\in M_d(\QQ)\) 可由 Hankel 矩阵 \(H_0=(a_{i+j})_{0\le i,j\le d-1}\) 与 \(H_1=(a_{i+j+1})_{0\le i,j\le d-1}\) 的闭式
\[
A=H_0^{-1}H_1=\frac{\mathrm{adj}(H_0)\,H_1}{\det(H_0)}
\]
直接构造。
当 \(a_n\in\ZZ\) 时 \(H_0,H_1\in M_d(\ZZ)\)，从而 \(A\) 的所有分母均整除 \(\det(H_0)\)。
\end{proposition}

\begin{remark}[分母下界与 \texorpdfstring{$p$}{p}-进刚性指数：Hankel 证书的算术刚性]\label{rem:pom-hankel-denominator-stiffness}
在 \(a_n\in\ZZ\) 的情形，上式表明 \(\det(H_0)\) 控制最小 realization 的有理分母。
对素数 \(p\) 可定义刚性指数
\[
\mathsf{Stiff}_p:=v_p\!\bigl(\det(H_0)\bigr),
\]
它给出在 \(\ZZ_p\) 中无歧义恢复 \(A\) 条目所需的最小 \(p\)-进精度门槛。
在本文的具体数据上，相应的 \(p\)-进刚性信号可由 Hankel 见证块行列式的赋值直接读出（注 \ref{rem:pom-hankel-stiffness-p5}）。
\end{remark}

\begin{remark}[非负整数碰撞核的搜索口径（建议作为新实验协议）]\label{rem:pom-nonneg-kernel-search-protocol}
命题 \ref{prop:pom-hankel-realization-protocol} 给出的 \(A\) 一般带有符号（或出现负条目），但它携带“最小线性实现”的确定结构。
为得到可解释的“有限状态碰撞核”\(A_q\)，可在 \(\mathrm{GL}_d(\ZZ)\) 上做可审计的相似变换搜索：
\[
A_q=P^{-1}AP,\qquad P\in\mathrm{GL}_d(\ZZ),
\]
目标是使 \(A_q\) 具有非负整数条目（从而可视为某个有限图/SFT 的邻接矩阵），并同步变换读出向量 \((\alpha,\beta)\) 得到同型的计数公式。
由于 \(P\) 为 unimodular，\(P^{-1}\in M_d(\ZZ)\)，该相似变换不会引入新的分母；若能找到使 \(A_q\in M_d(\ZZ_{\ge 0})\) 的 \(P\)，则该纯整基变换在保持最小维数的同时把分母完全消去（对比注 \ref{rem:pom-hankel-denominator-stiffness}）。
这一路径在 \(q=4\) 的五态核上已被脚本 \texttt{scripts/exp\_fold\_collision\_hankel\_realization\_q4.py} 以完全可复现方式实现；对更高阶 \(q\) 可按同一范式推广，并以“最小维数、最小最大条目、最小拓扑熵”等作为搜索准则，形成一个可无限上推的碰撞核谱系实验线路。
\end{remark}

\begin{corollary}[常数成本自由能求值器：由 Hankel 最小实现锁定 \texorpdfstring{$r_q$}{rq}]\label{cor:pom-hankel-free-energy-evaluator}
在命题 \ref{prop:pom-hankel-realization-protocol} 的设定下，记最小维数为 \(d\)。
则存在一个可审计偏移 \(\ell_q\ge 0\)（例如定义 \ref{def:pom-hankel-normal-form} 中使 Hankel 块可逆的最小偏移）使得仅由 \(2d\) 个样本点
\(
S_q(\ell_q+2),S_q(\ell_q+3),\dots,S_q(\ell_q+2d+1)
\)
即可构造 \(H_0^{(\ell_q)},H_1^{(\ell_q)}\)，从而得到最小实现矩阵 \(A=(H_0^{(\ell_q)})^{-1}H_1^{(\ell_q)}\)。
其特征值集合给出递推特征多项式的根，因此指数增长主尺度满足
\[
r_q=\rho(A),
\]
并且一旦找到某个非负整数 realization \(A_q=P^{-1}AP\)，则同样有 \(r_q=\rho(A_q)\)（Perron 根）。
从而对固定整数 \(q\)，自由能 \(\tau(q)=\log r_q\) 及“典型损失”母函数整点值
\(
\Lambda(q-1)=\log r_q-\log2
\)
（见命题 \ref{prop:pom-fiber-index-cgf} 与注 \ref{rem:pom-lambda-tau-shift}）都可在 \(O(d^3)\) 的矩阵代数成本内读出，而无需把 \(m\) 推至很大再拟合斜率。
\end{corollary}

\begin{definition}[Hankel-正规形：\texorpdfstring{$\mathrm{HANKELNF}_q$}{HANKELNFq}]\label{def:pom-hankel-normal-form}
固定 \(q\ge 2\)，令 \(a_n:=S_q(n+2)\)（\(n\ge 0\)），并令 \(d:=\mathrm{rank}(H)\) 为其 Hankel 秩（等价于最小递推阶数/最小线性实现维数）。
由于 \(\mathrm{rank}(H)=d\)，存在某个偏移 \(\ell\ge 0\) 使得 \(d\times d\) Hankel 块
\[
H_0^{(\ell)}:=(a_{\ell+i+j})_{0\le i,j\le d-1},
\qquad
H_1^{(\ell)}:=(a_{\ell+i+j+1})_{0\le i,j\le d-1},
\]
可逆。
令 \(\ell_q\) 为使 \(\det(H_0^{(\ell)})\neq 0\) 的最小偏移（一个可审计的整数证书；在实践中可由 Hankel-minor 搜索脚本输出）。
并定义 Hankel shift realization
\[
A:=(H_0^{(\ell_q)})^{-1}H_1^{(\ell_q)},
\qquad
\alpha^{\mathsf T}:=(a_{\ell_q},a_{\ell_q+1},\dots,a_{\ell_q+d-1}),
\qquad
\beta:=e_0.
\]
称
\[
\boxed{\;
\mathrm{HANKELNF}_q:=\bigl(d,\ell_q;\ H_0^{(\ell_q)},H_1^{(\ell_q)};\ A,\alpha,\beta\bigr)\;}
\]
为 \(q\)-moment 的\textbf{Hankel-正规形}：它是从有限段 \(a_{\ell_q},\dots,a_{\ell_q+2d-1}\) 直接、确定地构造出的最小 realization（命题 \ref{prop:pom-hankel-realization-protocol}）。
\end{definition}

\begin{proposition}[moment 的有限采样判等：\texorpdfstring{$2d$}{2d} 点决定全序列]\label{prop:pom-moment-finite-sampling-decidable}
在定义 \ref{def:pom-hankel-normal-form} 的设定下，设两条实现路径（例如两种编译/状态编码）给出两条 \(q\)-moment 序列 \(S_q,S'_q\)，并且二者的 Hankel 秩同为 \(d\)。
若
\[
S_q(m)=S'_q(m)\qquad(m=\ell_q+2,\ \ell_q+3,\ \dots,\ \ell_q+2d+1),
\]
则它们具有同一个 Hankel-正规形 \(\mathrm{HANKELNF}_q\)，从而对所有 \(m\ge 2\) 都有 \(S_q(m)=S'_q(m)\)；并且由该正规形导出的全部指纹（递推、特征多项式、Perron 根 \(r_q\)、以及后续的压力/预算读出等）完全一致。
\end{proposition}
\begin{proof}
将条件改写为 \(a_n=a'_n\) 对所有 \(\ell_q\le n\le \ell_q+2d-1\) 成立。
因此两侧由该长度 \(2d\) 的窗口构造出的 \(H_0^{(\ell_q)},H_1^{(\ell_q)}\)（以及 \(A=(H_0^{(\ell_q)})^{-1}H_1^{(\ell_q)}\)、\(\alpha,\beta\)）逐项相同，故 \(\mathrm{HANKELNF}_q=\mathrm{HANKELNF}'_q\)。
由 \(a_n=\alpha^{\mathsf T}A^n\beta\) 的最小 realization 表达（命题 \ref{prop:pom-hankel-realization-protocol}）立得对所有 \(n\ge 0\) 有 \(a_n=a'_n\)，即对所有 \(m\ge 2\) 有 \(S_q(m)=S'_q(m)\)。
其余指纹一致性由同一正规形共享同一 \(A\) 与同一特征多项式直接推出。证毕。
\end{proof}

\begin{remark}[双路一致性审计：residue-DP 拟合 vs.\ realization Perron 根]\label{rem:pom-dp-vs-realization-audit}
对每个固定 \(q\)，可并行输出两条可复核读出：
\begin{itemize}
  \item residue-DP 在有限窗口上计算 \(S_q(m)\) 并据此拟合得到的 \(\widehat r_q\)；
  \item 由递推/伴随矩阵或 Hankel 最小实现（以及其非负整数核 \(A_q\)）读出的精确 \(r_q=\rho(A_q)\)。
\end{itemize}
二者应一致；若不一致，则要么 DP 窗口不足、要么 realization 并非同一序列的实现（或存在边界项未吸收进起止向量）。
\end{remark}

\subsubsection{内生 Gibbs 权重：由纤维体积诱导的输出平衡态与 transfer matrix}\label{subsubsec:pom-gibbs-pi}
本文已将折叠输出分布定义为
\[
\pi_m(x):=\frac{d_m(x)}{2^m}\qquad(x\in X_m),
\]
其中 \(d_m(x)=|\Fold_m^{-1}(x)|\) 是纤维体积。
这一形式可以被理解为一个完全内生的 Gibbs 权重：令能量
\[
E_m(x):=-\log d_m(x),
\]
则
\[
\pi_m(x)=\frac{e^{-E_m(x)}}{Z_m},
\qquad
Z_m:=\sum_{x\in X_m}e^{-E_m(x)}=\sum_{x\in X_m}d_m(x)=|\Omega_m|=2^m.
\]
因此 \(\pi_m\) 不是外加的基准分布，而是折叠映射作为“粗粒化”在稳定语法空间上诱导的自然平衡态：纤维越大（预像自由度越高），其输出质量越大。
\par
进一步，若希望把该 Gibbs 权重提升为一个可计算的 transfer matrix（以替代仅用 Parry/PF 作为外部基线的口径），一个可审计路线是：从 \(\pi_m\) 的加权块频率出发，构造一个有限阶 Markov（等价地：有限程 Gibbs）近似，并以 Perron--Frobenius 抽取其稳定态。

\begin{remark}[可复现协议：由 \texorpdfstring{$\pi_m$}{pi_m} 拟合有限阶 transfer matrix]\label{rem:pom-gibbs-markov-fit}
固定 \(m\)（取中等分辨率，如 \(m=18\)），枚举得到 \(\pi_m\) 后，对某个小块长 \(\ell\ge 2\) 计算 \(\pi_m\)-加权的 \(\ell\)-块频率，
并由最大似然得到一个 \((\ell-1)\)-阶 Markov 链（其状态为空间 \(X_{\ell-1}\) 的合法块）。
该链即为一个有限状态 transfer matrix 的概率规范化版本；其诱导的 1-步转移可与 Parry 基线作定量对照，从而得到“内生势函数偏离最大熵基线”的可审计证书。
表 \ref{tab:fold_output_gibbs_markov_fit} 给出一个代表性输出（脚本 \texttt{scripts/exp\_fold\_output\_gibbs\_markov\_fit.py} 一键生成）。
\end{remark}
\input{sections/generated/tab_fold_output_gibbs_markov_fit}

\begin{remark}[熵率/KL 证书：内生 Gibbs 权重与 Parry 最大熵基线的微比特级差异]\label{rem:pom-gibbs-parry-entropy-gap}
在 golden-mean SFT 的两状态邻接
\(
A=\begin{psmallmatrix}1&1\\[2pt]1&0\end{psmallmatrix}
\)
下，Parry（PF）基线 \(P_{\mathrm{Parry}}\) 具有最大熵率
\(
h(P_{\mathrm{Parry}})=\log\varphi
\)
（其中 \(\varphi\) 为 Perron 根）。
对任意与之同支撑的两状态 Markov 规格 \(P\)，其相对 Parry 的 KL 散度率满足
\[
D_{\mathrm{KL}}^{\mathrm{rate}}(P\|P_{\mathrm{Parry}})=\log\varphi-h(P),
\]
从而“最大熵缺口”可等价输出为一份 KL-率证书。
\par
对表 \ref{tab:fold_output_gibbs_markov_fit} 的拟合输出（\(m=18,\ell=4\)），该缺口已降至微比特级；脚本 \texttt{scripts/exp\_fold\_output\_gibbs\_markov\_fit.py} 同步输出如下可复算数值证书：
\input{sections/generated/eq_fold_output_gibbs_entropy_gap}
该量级支持“折叠诱导的内生势函数偏离最大熵基线主要体现为常数级指纹，而非熵率层面的结构崩塌”的解释框架。
\end{remark}

\begin{remark}[接口闭环：碰撞谱 \texorpdfstring{$\leftrightarrow$}{<->} 典型空间损失的三条可审计推论]\label{rem:pom-interface-loop-collision-loss}
本节把两条接口（Hankel 最小实现与内生 Gibbs 权重）并排写出后，可得到如下三条闭环推论（均由文内机制直接推出，并可用现有脚本核验）：
\begin{itemize}
  \item \textbf{（同构）}对 \(X\sim\pi_m\) 与 \(L_m=\log d_m(X)\)，有
  \(\mathbb{E}_{\pi_m}[e^{(q-1)L_m}]=S_q(m)/2^m\)，从而 \(\Lambda(t)=\tau(t+1)-\log2\)，并给出 \(\mu=\tau'(1)\)、\(\sigma^2=\tau''(1)\)（命题 \ref{prop:pom-fiber-index-cgf} 与注 \ref{rem:pom-lambda-tau-shift}）。
  \item \textbf{（零参数夹逼）}\(\tau\) 的凸性把 \(\mu\) 用三点端点夹死：\(\log(2/\varphi)\le \mu\le \log(r_2/2)\)；并且有限 \(m\) 上已给出逐项可审计夹逼（推论 \ref{cor:pom-mu-secant-audit}）。
  \item \textbf{（常数成本求值器）}一旦由 Hankel 口径恢复出某个 \(q\) 的最小 realization（并进一步找到非负整数核 \(A_q\)），则 \(r_q=\rho(A_q)\) 立刻锁定，从而同时锁定 \(\tau(q)=\log r_q\) 与 \(\Lambda(q-1)=\log r_q-\log2\)，并可与 residue-DP 拟合输出做双路一致性审计（推论 \ref{cor:pom-hankel-free-energy-evaluator} 与注 \ref{rem:pom-dp-vs-realization-audit}）。
\end{itemize}
\end{remark}

